\newcommand{\ModuleID}{EL-P1}
\newcommand{\ModuleIDKey}{P1}
\newcommand{\ModuleTitle}{Textual Prosody and Quantity}
\newcommand{\ModuleStage}{Advanced Poetry Branch}
\newcommand{\ModuleUnit}{P1}
\newcommand{\ModuleOutcome}{Read quantitative verse from the printed page by identifying syllables, quantity, elision, feet, and line structure without turning the course into pronunciation instruction}
\newcommand{\ModulePrerequisites}{The complete common trunk through EL-M22, especially reliable morphology, poetic word order, textual annotation, and controlled composition}
\newcommand{\ModuleNext}{EL-P2, Hexameter}
\newcommand{\ModuleBeforeContinuing}{Restore the supplied six-foot senarius frame and its bounded substitutions; from a printed line, divide words into syllables, distinguish quantity by nature and by position, mark defensible elisions, test rather than guess the metrical pattern, preserve textual uncertainties, and explain every mark without relying on a pronunciation guide.}
\newcommand{\ModuleDelayNote}{}
\newcommand{\ModuleExtensionPractice}{Return to the exact memory set, reading, scansion, commentary, composition task, or worksheet that exposes the uncertainty; then use different verse for the later return.}
\newcommand{\ModuleLessonSource}{\CourseRoot/shared/content/poetry-advanced/all}
\newcommand{\ModuleExerciseSource}{\CourseRoot/shared/exercises/poetry/all}
\newcommand{\ModuleWorkedBank}{P}
\newcommand{\ModuleExerciseBank}{P}
\newcommand{\ModuleWorksheetVariants}{A,B,C,D}
\newcommand{\ModuleMemorySet}{P1}
\newcommand{\ModuleReaderSet}{P1}
\newcommand{\ModuleScopeNote}{This branch is page-based textual prosody. The packet prints both the Phaedrus fable window and the verified prologue wording that names senarian verse; it neither supplies nor presupposes a spoken pronunciation system.}
\newcommand{\ModulePractice}{%
  \begingroup
  \SelectCourseUnits{P1}
  % Integrated practice for the advanced poetry branch. Each selected block is
% rendered once as fresh work and again as models. R identifies received text
% already documented in the branch source records; P identifies project-made
% instructional Latin or analysis.

\begin{courseunit}{P1}
\SubmoduleExtension
  {Quantity-evidence sorting}
  {A metrical mark becomes dependable only when its reason remains attached
  to it. Natural quantity, positional length, elision, and final anceps answer
  different questions.}
  {Make four columns headed \emph{nature}, \emph{position}, \emph{elision
  environment}, and \emph{not yet established}. Sort and explain these P
  examples: \Latin{caelum, aurum, poena, terra, multus, patrem, multam
  amicitiam}. Do not assign a metrical foot.}
  {Write the three ordinary diphthongs, the positional rule, the elision
  environment, and the meaning of final anceps.}
  {\Latin{caelum, aurum, poena} display \Latin{ae, au, oe}, ordinarily long
  by nature. In \Latin{terra} and \Latin{multus}, the preceding syllable is
  long by position before two consonants. \Latin{patrem} presents a stop-plus-
  liquid environment that should be marked for testing, not silently forced
  into one treatment. \Latin{multam amicitiam} presents final vowel plus
  \Latin{-m} before an initial vowel and therefore an ordinary elision
  environment. No item by itself establishes a foot.}
\SubmoduleExtension
  {Grammar before scheme}
  {Verse still has clause structure. Recovering that structure first prevents
  a proposed meter from manufacturing forms or relations.}
  {Using R Phaedrus 1.1.1--2, write the two finite-verb spines; attach
  \Latin{ad rivum eundem}, \Latin{siti compulsi}, and \Latin{superior}; then
  make a P prose-order reconstruction. Only after that begin a separate
  syllable-and-quantity ledger.}
  {Parse \Latin{eundem, venerant, siti, compulsi, superior}.}
  {The spines are \Latin{lupus et agnus venerant} and \Latin{lupus stabat}.
  One P reconstruction is \Latin{Lupus et agnus, siti compulsi, ad eundem
  rivum venerant. Lupus superior stabat.} \Latin{Eundem} is accusative
  masculine singular with \Latin{rivum}; \Latin{venerant} is third-person
  plural pluperfect active indicative; \Latin{siti} is ablative singular of
  cause; \Latin{compulsi} is nominative masculine plural; \Latin{superior} is
  nominative masculine singular and predicative of the wolf.}
\SubmoduleExtension
  {A reproducible scansion ledger}
  {Repeated work should expose uncertainty rather than conceal it. A second
  return is useful when it can reproduce the first return's reasons without
  copying its marks.}
  {Choose either of the two R Phaedrus lines printed at the opening of P1,
  preferably the line not used for your first ledger. Copy its received
  wording, divide it into syllables, and make columns for proposed quantity,
  reason, possible elision, metrical position, and uncertainty. Close the
  packet, rebuild the ledger, and reconcile only the differences for which
  you can state evidence.}
  {State the pure six-foot senarius frame, its bounded substitutions, and the
  order of operations from source transcription through scheme testing.}
  {The pure frame has six short--long iambic feet, with final anceps; a
  tribrach may replace an iamb outside the last foot, while odd feet may also
  receive the stated spondee, dactyl, or anapaest substitutions. A sound
  record preserves the source text; divides syllables; marks secure
  natural quantity; adds positional environments; notes possible elision;
  leaves unresolved evidence open; and only then tests a stated scheme. A
  later mismatch is repaired by checking evidence, not by silently changing a
  lexical quantity or the received wording.}
\end{courseunit}

\begin{courseunit}{P2}
\SubmoduleExtension
  {Foot boundary and word boundary}
  {Hexameter is organized in feet, but syntax and words do not begin and end
  at every foot boundary. Keeping two rows prevents a visual division from
  becoming a false grammatical division.}
  {Copy R \Latin{Arma virumque cano, Troiae qui primus ab oris}. On one row
  mark its six metrical divisions from the teaching; on a second row mark word
  boundaries; on a third row write the clause relations. Circle every place
  where the three rows do not coincide.}
  {Draw the hexameter frame: the alternatives in feet one through four, the
  normal fifth foot, and the final foot.}
  {The taught metrical grouping is \Latin{arma vi | rumque ca | no Tro |
  iae qui | primus ab | oris}. The words remain \Latin{arma / virumque / cano
  / Troiae / qui / primus / ab / oris}; \Latin{arma} and \Latin{virum} are
  objects of \Latin{cano}, while \Latin{qui} begins a relative clause whose
  finite verb is delayed. Feet one through four permit dactyl or spondee, the
  fifth is normally a dactyl, and the sixth is long plus anceps.}
\SubmoduleExtension
  {One epic sentence}
  {The first four lines teach sustained reading because their finite and non-
  finite forms distribute one proposition across several verse boundaries.}
  {Make a sentence map for R \textit{Aeneid} 1.1--4. Begin with \Latin{cano}
  and its objects; nest the relative clause under \Latin{virum}; place
  \Latin{venit}; then attach origin, destinations, modifiers, and the four
  expressions of circumstance or cause. Write a P prose-order aid underneath.}
  {Recover the functions of \Latin{Troiae, Italiam, fato, Lavina litora,
  terris, alto, vi superum, ob iram}.}
  {A defensible map begins \Latin{cano arma et virum}; \Latin{qui} refers to
  \Latin{virum}, and \Latin{venit} governs the relative clause. \Latin{ab oris
  Troiae} gives origin; \Latin{Italiam} and \Latin{Lavina litora} are
  destinations; \Latin{fato} accompanies \Latin{profugus}; \Latin{terris} and
  substantival \Latin{alto} mark land and sea; \Latin{vi superum} and
  \Latin{ob iram} state forces or causes. A P prose order may follow the model
  in the teaching, but it is not a Virgilian variant.}
\SubmoduleExtension
  {Prose plan before hexameter}
  {Composition develops from an exact proposition. Meter may rearrange or
  refine sound Latin, but it cannot excuse a lost subject, false quantity, or
  altered destination.}
  {Begin with P \Latin{Peregrinus, spe impulsus, per tempestates ad urbem
  sanctam venit.} Make a syntax ledger and a quantity ledger. Try one
  hexameter line, record every substitution, and return to the prose plan if a
  sound line cannot yet be made.}
  {Without notes, list the relations that the verse must preserve.}
  {The pilgrim remains the subject; \Latin{spe} supplies the motivating cause;
  \Latin{impulsus} agrees with the subject; \Latin{per tempestates} marks the
  traversed circumstance; \Latin{ad urbem sanctam} remains the destination;
  and \Latin{venit} gives completed arrival. A documented prose result is
  sounder than a line purchased with broken grammar or unsupported quantity.}
\end{courseunit}

\begin{courseunit}{P3}
\SubmoduleExtension
  {The two-line contract}
  {An elegiac couplet is not two interchangeable hexameters. Its second line
  has a fixed central division and a different closure.}
  {Draw from memory a six-foot hexameter above a divided pentameter. Label the
  two halves of the pentameter and its closing pattern. Then use R Ovid,
  \textit{Tristia} 1.1.1--2, to mark line form, clause boundary, and semantic
  turn on three separate rows.}
  {State what normally happens to a word at the pentameter's central division.}
  {The first line is a dactylic hexameter; the second is a pentameter made of
  two balanced halves around a fixed division. A word normally does not
  bridge that division. The formal boundary is visible from the verse form,
  while grammar and contrast establish the semantic turn.}
\SubmoduleExtension
  {Address, person, and exclusion}
  {Ovid's couplet changes grammatical person as the speaker distinguishes the
  book's freedom from his own condition.}
  {Color-code second-person address, first-person reference, and third-person
  renaming in \Latin{Parve---nec invideo---sine me, liber, ibis in urbem: ei
  mihi, quod domino non licet ire tuo!} Then write a P prose order and explain
  the construction with \Latin{licet}.}
  {Parse \Latin{parve, liber, ibis, me, mihi, domino tuo, ire}.}
  {\Latin{Parve} and \Latin{liber} are vocatives addressing the book;
  \Latin{ibis} is second-person singular future active indicative;
  \Latin{me} and \Latin{mihi} refer to the speaker; \Latin{domino tuo} is
  dative with impersonal \Latin{licet}; and \Latin{ire} completes it. One P
  order is \Latin{Parve liber, sine me in urbem ibis---nec invideo. Ei mihi,
  quod domino tuo ire non licet!}}
\SubmoduleExtension
  {Designing a turn before making verse}
  {A semantic turn can be planned and tested in prose before metrical choices
  make revision more difficult.}
  {Write two P sentences: a traveler reaches a church, but remembers an absent
  friend. Divide the thought into opening movement and turn; verify every
  form; then attempt an elegiac couplet with separate quantity, syntax, and
  revision ledgers.}
  {Name six reasons for returning a metrical attempt to prose.}
  {A compact P plan is \Latin{Viator tandem ad ecclesiam pervenit. Laetatur,
  sed amici absentis memor est.} Return to prose if meter demands false
  quantity, wrong case, broken agreement, an unintended proposition,
  indefensible order, or undocumented borrowing. The plan remains useful even
  when no sound couplet has yet emerged.}
\end{courseunit}

\begin{courseunit}{P4}
\SubmoduleExtension
  {One lyric contract repeated}
  {A fixed lyric measure is tested by recurrence. One difficult line does not
  authorize a newly invented scheme.}
  {Draw eight rows for R Horace, \textit{Carmina} 1.11. Put the same greater-
  Asclepiadean template beside each row: opening long--long base, three
  long--short--short--long choriambic sequences, and the final
  short--anceps close, with the two stable internal divisions. Test two lines
  now and leave every unresolved quantity visible.}
  {Define a choriamb and explain why punctuation cannot establish a metrical
  division by itself.}
  {A choriamb is long--short--short--long. Each line is tested against the same
  declared architecture. Punctuation records syntactic or editorial
  articulation; a metrical division belongs to the scheme. The two can
  coincide, but one does not prove the other.}
\SubmoduleExtension
  {Counsel through a long syntax}
  {The famous final imperative receives its force from a sequence of
  prohibition, alternatives, endurance, and concentrated counsel.}
  {Make a finite-verb ledger for the ode. Group \Latin{quaesieris, dederint,
  temptaris}; then \Latin{erit, tribuit, debilitat}; then \Latin{sapias,
  liques, reseces}; and finally \Latin{loquimur, fugerit, carpe}. State the
  clause or discourse work of each group.}
  {Reconstruct the ode's argument without beginning from \Latin{carpe diem}.}
  {The first group prohibits inquiry and embeds what the gods have allotted;
  the second opens unknown alternatives and the present winter; the third
  concentrates practical counsel; the fourth opposes present speech to time
  already fled and closes in imperative force. \Latin{Carpe diem} is the
  culmination of that syntax, not a detachable summary of the entire poem.}
\SubmoduleExtension
  {A controlled project comparison}
  {Similar brevity or devotional force does not establish common meter,
  authorship, or historical dependence. A P comparison can still isolate
  grammar and structure.}
  {Compare Horace with the project-built hymn-style stanza \Latin{Christe,
  lumen cordium, mentes fessas recrea; da quietem labori, dirige nos ad
  patriam.} Make separate rows for source class, clause density, commands,
  balance, line closure, and any claimed pattern.}
  {Parse the P stanza and state what its label prevents you from claiming.}
  {\Latin{Christe} is vocative; \Latin{lumen cordium} is apposition;
  \Latin{mentes fessas} is object of \Latin{recrea}; \Latin{quietem} is
  object and \Latin{labori} dative with \Latin{da}; \Latin{nos} is object of
  \Latin{dirige}; \Latin{ad patriam} gives direction. The P label prevents an
  invented ancient author, date, meter, liturgical use, or witness.}
\end{courseunit}

\begin{courseunit}{P5}
\SubmoduleExtension
  {Five claims kept separate}
  {A Missal witness establishes printed wording and liturgical placement in
  that edition. Authorship, date, literary form, and a proposed scriptural
  relationship require their own evidence.}
  {For the Pentecost Sequence make a claim ledger with rows for witness,
  liturgical use, authorship, date, poetic organization, and scriptural source.
  Fill only what the supplied source record establishes; mark the evidence
  still needed for every further claim.}
  {Without notes, distinguish a received witness fact, grammatical analysis,
  source-grounded inference, and unverified lead.}
  {MR62-FORM-PENT establishes the wording and Pentecost placement printed in
  the identified 1962 Missal. It does not by itself establish original
  authorship, first date, earlier transmission, complete metrical history, or
  dependence on a particular biblical passage. Grammar can be analyzed from
  the text; broader claims remain separately sourced or explicitly open.}
\SubmoduleExtension
  {Apposition and imperative triads}
  {The sequence moves from naming the addressee and the relief desired to
  direct petitions for restoration.}
  {Map \Latin{Consolator optime, dulcis hospes animæ, dulce refrigerium} and
  \Latin{In labore requies, in æstu temperies, in fletu solacium}. Then make
  a syntax table for \Latin{lava, riga, sana / flecte, fove, rege} with each
  accompanying free relative.}
  {Reconstruct the six imperative-plus-\Latin{quod} clauses in order.}
  {The opening names one addressee through vocative apposition, followed by
  three balanced nominal descriptions of relief. Each of the six imperatives
  governs a neuter free relative whose \Latin{quod} is subject of
  \Latin{est}; its predicate adjective identifies what is dirty, dry,
  wounded, rigid, cold, or wandering. The two triads move from cleansing and
  healing to correction, warming, and guidance.}
\SubmoduleExtension
  {One structure, newly composed}
  {Imitation is most transparent when it borrows one declared structure and
  records which words are new.}
  {Label as P and parse \Latin{In via comes, in nocte lumen, in periculo
  auxilium. Erige quod est lapsum, confirma quod est infirmum, illumina quod
  est obscurum.} Revise one lexical field and one imbalance, then record the
  sequence structure imitated and every exact received phrase retained.}
  {Cover the model and rebuild its three nominal members and three imperative
  clauses.}
  {The first triad uses prepositional settings plus nominative predicate or
  appositional nouns. In the second, each singular imperative governs a free
  relative: neuter \Latin{quod} is subject of \Latin{est}, and
  \Latin{lapsum, infirmum, obscurum} are predicate adjectives. The source
  ledger should identify structural imitation of the Pentecost Sequence and
  must not present this P wording as received or liturgical.}
\end{courseunit}

\begin{courseunit}{P6}
\SubmoduleExtension
  {Three independent records}
  {Quantity, page-based recurrence, and visible line-ending
  correspondence can coexist in one text without becoming the same kind of
  evidence.}
  {For \Latin{In labore requies, in æstu temperies, in fletu solacium}, make
  one grammar ledger, one syllable-count and parallel-structure ledger, one
  quantity ledger, and one written-ending comparison. Give each conclusion
  the name of the ledger that supports it.}
  {Define quantitative structure and page-based textual pattern; then explain
  why visible endings alone do not prove rhyme or assonance.}
  {Grammar records the three prepositional circumstances and balanced nouns.
  The pattern ledger records recurring counts, parallel grammar, and line
  grouping; the quantity ledger records natural and positional evidence plus
  uncertainty; the ending ledger compares visible material such as
  \Latin{-quies, -temperies, -solacium}. A sound classification as rhyme or
  assonance would additionally need phonological evidence not supplied here.}
\SubmoduleExtension
  {Imperatives as a memory architecture}
  {A finite-form spine can preserve the large movement of a poem before every
  lexical detail has become secure.}
  {From memory restore \Latin{veni---emitte---veni---veni---veni---reple---
  lava---riga---sana---flecte---fove---rege---da---da---da---da}. Group the
  verbs under arrival, illumination, restoration, guidance, and gift. Then
  compare the chain with the printed witness and record every omission or
  transposition.}
  {Rebuild the chain twice: once by form and once by semantic group.}
  {Arrival uses four occurrences of \Latin{veni}; sending or illumination and
  filling use \Latin{emitte, reple}; restoration uses \Latin{lava, riga,
  sana}; correction and guidance use \Latin{flecte, fove, rege}; gift uses
  four occurrences of \Latin{da}. The later comparison should preserve
  differences as useful retrieval evidence rather than silently erase them.}
\SubmoduleExtension
  {Patterned imitation with a grammar audit}
  {Visible recurrence is worth keeping only when the repeated lines also have
  sound case, agreement, government, and meaning.}
  {Analyze the project-built hymn-style stanza \Latin{Veni, lux fidelium;
  mentes nostras illumina. In labore da solacium; ad bonum nos dirige.} Then
  write four to eight P lines with one repeated grammatical pattern and two
  exact visible ending correspondences. Keep a prose plan and revision list.}
  {List what may be revised for pattern and what may not be falsified.}
  {\Latin{Veni, illumina, da, dirige} are singular imperatives;
  \Latin{lux fidelium} is vocative plus genitive; \Latin{mentes nostras},
  \Latin{solacium}, and \Latin{nos} are objects; \Latin{in labore} gives a
  circumstance and \Latin{ad bonum} direction. Syllable count, order,
  parallel syntax, and ending vocabulary may be revised; grammar, meaning,
  source labels, and asserted quantity may not be falsified.}
\end{courseunit}

\begin{courseunit}{P7}
\SubmoduleExtension
  {The seven-layer commentary from memory}
  {A fixed order keeps witness fact, analysis, and interpretation from
  silently trading places.}
  {Draw seven boxes and fill them, without notes, with: witness and boundary;
  text, normalization, and any supplied material variants; meter or visible
  pattern, with quantity where applicable; syntax and P prose order;
  diction and rhetoric; supported source or allusion context; interpretation,
  limits, and P translation. Then audit one Phaedrus claim through all seven.}
  {Reproduce the seven layers in order and give one question each answers.}
  {The sequence fixes what text is under discussion before analysis, preserves
  textual uncertainty before scansion, grounds syntax before rhetoric, and
  places interpretation only after the supporting evidence. A source-context
  claim needs an actual source window; unresolved quantity reappears among the
  stated limits when it affects the conclusion.}
\SubmoduleExtension
  {From exact form to proportionate thesis}
  {Phaedrus first gives shared need and location, then unequal position, then
  characterizes the aggressor before his accusation.}
  {Using R Phaedrus 1.1.1--6, make an observation table for
  \Latin{eundem, siti compulsi, superior, longeque inferior, fauce improba,
  latro, iurgii causam intulit}. From the table write one thesis, one plausible
  counter-observation, and one explicit limit.}
  {State the difference between textual fact and source-grounded inference.}
  {A proportionate P thesis is: \English{Phaedrus turns physical elevation
  into an image of unequal power before the wolf manufactures a verbal
  grievance.} The exact spatial adjectives, \Latin{latro}, and the narrative
  sequence support it. The thesis is an inference, not source wording; the
  excerpt alone does not establish a universal theory of fable, the poet's
  biography, or a specific political allegory.}
\SubmoduleExtension
  {A compact commentary dossier}
  {Repeated short dossiers make source discipline habitual before a longer
  commentary demands sustained argument.}
  {Choose four to twelve lines from one of the verified Virgil, Horace, or
  Pentecost texts reprinted in P7. Supply the stable source ID and locus,
  normalization note, metrical or visible-pattern ledger suited to the text,
  syntax map, one rhetoric claim, one verified context window if used, a P
  translation, a thesis, counterevidence, and a limit. Then reduce the dossier
  to 150 words without dropping a claim class.}
  {Close the packet and list every item that must remain when a commentary is
  shortened.}
  {A sound short dossier still names the witness and boundary, separates text
  from normalization, gives reasons rather than unexplained metrical marks,
  accounts for the finite spine and significant non-finite forms, grounds its
  rhetorical claim in exact Latin, labels its translation P, and states what
  the evidence does not establish. Unsupported source context is omitted, not
  implied.}
\end{courseunit}

\begin{courseunit}{P8}
\SubmoduleExtension
  {The prose plan controls the poem}
  {A declared thought and form make later substitutions auditable. The poem
  should not discover its proposition accidentally while chasing a pattern.}
  {Choose one bounded form from the packet. Write speaker, addressee, principal
  proposition, grammatical spine, movement of thought, register, and essential
  images in prose. Divide that unchanged prose into line-sized movements before
  making any metrical or rhythmic substitution.}
  {Name the four bounded form choices and the fields of the prose plan.}
  {The choices are four hexameters, two elegiac couplets, one Horatian line
  pattern sustained for four lines, or a six- to twelve-line rhythmic hymn-
  style stanza. The plan fixes speaker and addressee, proposition, grammar,
  sequence, images or petitions, register, and relation between movements.}
\SubmoduleExtension
  {Four ledgers, one draft}
  {Quantity, syntax, revision, and source records inspect different kinds of
  integrity. A line is not sound merely because one ledger is complete.}
  {Audit one drafted line in four tables: syllable, quantity, reason, and
  metrical place; word, morphology, governor, and function; old wording, new
  wording, and reason; feature or phrase, source ID and locus, and kind of use.
  Give the line a verdict of \emph{retain}, \emph{revise}, or \emph{return to
  prose}.}
  {Rebuild the columns of all four ledgers without notes.}
  {The quantity ledger exposes unsupported length or resolution; the syntax
  ledger exposes unattached or false forms; the revision ledger preserves
  semantic and structural costs; the source ledger distinguishes quotation,
  adaptation, structural model, echo, and general influence. The verdict must
  cite exact evidence from all relevant ledgers.}
\SubmoduleExtension
  {Final documentation and fresh return}
  {Original verse remains clearly original when the final presentation makes
  its pattern, sources, translation, changes, and uncertainty inspectable.}
  {Prepare a P-labeled poem with declared form, prose statement of thought,
  pattern record, P translation, syntax notes, revision summary, source
  ledger, remaining uncertainty, and P7 commentary. Put it aside; on a later
  return, rescan and construe it from a clean copy before reading the ledgers.}
  {Translate and parse \Latin{Nihil tamen auctori antiquo attribuo quod ipse
  composui}.}
  {\English{Yet I attribute to no ancient author what I myself composed.}
  \Latin{Nihil} is the object of \Latin{attribuo}; \Latin{auctori antiquo} is
  dative; \Latin{quod} is the object of \Latin{composui} and refers to what
  was composed; \Latin{ipse} emphasizes the composer. A fresh return records
  defects honestly and may end with a documented unfinished poem rather than
  a falsely secure one.}
\end{courseunit}

  \endgroup}
